

\section{Pfaffian wall normal form on type-II walls}
\label{app:pfaffian-wall-normal-form}

The three local Pfaffian wall computations supply the
sign data used in Chapter \ref{ch:pfaffian-dirac} for the
orientation-character match
\[
\epsilon_o = \nu_{\Delta_5}\quad\text{on}\quad
W^{(2)}(\La^{2,1}_{II}),
\]
after the global orientation transport of
Theorems~\ref{thm:G-O1}, \ref{thm:G-O1-plus}, and
\ref{thm:G-O2-orbit-transport}. Every numerical
identity of \S\ref{app:pfaffian-wall-three-walls} and
\S\ref{app:pfaffian-wall-recognition} is independently verified by
exact rational lattice arithmetic and exact rational expansion of the
Borcherds product.

\subsection{The three simple type-II walls}
\label{app:pfaffian-wall-three-walls}

The hyperbolic root lattice of signature \((2,1)\) is
\(\La^{2,1} = \HypPlan \oplus \langle 2 \rangle\), with basis
\((f_2, f_3, f_{-2})\) and bilinear form
\((f_2, f_{-2}) = -1\), \((f_3, f_3) = 2\). Its index-four type-II
sublattice is
\[
\La^{2,1}_{II}
= \{m f_2 + l f_3 + n f_{-2} \in \La^{2,1} \mid
m \equiv n \equiv 0 \pmod 2\},
\]
abstractly \(\La^{2,1}_{II}\simeq U(4)\oplus\langle 2\rangle\) with
discriminant \(-32\); the reflection-group index
\([W^{(2)}(\La^{2,1}):W^{(2)}(\La^{2,1}_{II})]=6\) tracks the orbit
data on the chamber-walls and is independent of the sublattice index.
The three simple type-II roots are
\[
\delta_1 = 2 f_2 - f_3, \qquad
\delta_2 = 2 f_{-2} - f_3, \qquad
\delta_3 = f_3,
\]
with pairwise pairings forming the type-II Cartan matrix
\[
A = \bigl((\delta_i, \delta_j)\bigr)_{1 \le i, j \le 3}
=
\begin{pmatrix}
2 & -2 & -2 \\
-2 & 2 & -2 \\
-2 & -2 & 2
\end{pmatrix},
\qquad
\det A = -32.
\]
The matrix is symmetric generalized Cartan, hyperbolic (one negative
eigenvalue, two positive), and the type-II Weyl group
\(W^{(2)}(\La^{2,1}_{II})\) generated by
\(s_{\delta_1}, s_{\delta_2}, s_{\delta_3}\) is the Coxeter group with
relations \(s_{\delta_i}^2=1\) and no braid relation between distinct
generators; the off-diagonal value \(-2\) is the infinite Coxeter
exponent.  It acts on the hyperbolic plane \(\R\Poly_{II}\). The Weyl
vector \(\rho\) determined
by \((\rho, \delta_i) = -1\) for \(i = 1, 2, 3\) is
\[
\rho = \tfrac{1}{2}(\delta_1 + \delta_2 + \delta_3)
     = f_2 - \tfrac{1}{2} f_3 + f_{-2}.
\]
The Cartan matrix, the Weyl vector, the reflection identity
\(s_{\delta_i}(\delta_i) = -\delta_i\), and the preservation of the
form under \(s_{\delta_i}\) are all direct computations in the BKM
Gram.

\begin{remark}[Type-II versus type-I]
The lattice \(\La^{2,1}\) admits both type-I roots
(divisibility \((\delta, \La^{2,1}) = \Z\), simple system
\(\{f_{-2} - f_2, f_2 - f_3, f_{-2} - f_3\}\), trivial multiplier) and
type-II roots (divisibility \(2\Z\), simple system
\(\{\delta_1, \delta_2, \delta_3\}\) above, multiplier \(-1\) on each
generator). The Borcherds--Gritsenko--Nikulin denominator algebra
\(\fg_{\Delta_5}\) is the type-II algebra, with denominator
\(64^{-1} \Delta_5(2Z)\) on \(\La^{2,1}_{II}\); the doubling \(2Z\) is
the passage from the type-I to the type-II lattice, and the constant
\(64 = 2^6\) is the theta-normalization
\([q^{1/2} r^{1/2} s^{1/2}] \Delta_5\). The Pfaffian wall calculation
is type-II throughout.
\end{remark}

\subsection{The rank-one type-II Pfaffian crossing}
\label{app:pfaffian-wall-rank-one}

For each simple type-II wall
\(\delta \in \{\delta_1, \delta_2, \delta_3\}\), the retained
rank-one chart in the datum \((O2_{\mathrm{atlas}})\) has local model
\[
U_\delta = T_\delta \times (-\varepsilon, \varepsilon),
\qquad
s_\delta(t, u) = (t, -u),
\]
with wall locus \(T_\delta \times \{0\}\) coinciding with the fixed
locus of \(s_\delta\).  The reduced normal self-Ext factor on this
retained chart is the rank-one odd Pfaffian crossing
\[
K^{\mathrm{cross}}_\delta = [\R \xrightarrow{\,u\,} \R],
\qquad
D_u = \begin{pmatrix} 0 & u \\ -u & 0 \end{pmatrix},
\qquad
\mathrm{Pf}(D_u) = u.
\]
Choose an oriented tangent Pfaffian line
\(\mathscr L^{\mathrm{tan}}_\delta\) on \(T_\delta\) with
\(s_\delta\)-invariant unit \(\upsilon_\delta\):
\[
s_\delta(\upsilon_\delta) = \upsilon_\delta.
\]
The total local Pfaffian section is
\[
\mathrm{pf}_\delta = \upsilon_\delta\, u.
\]

\begin{proposition}[Local Pfaffian sign; conditional on the retained rank-one wall chart]
\label{prop:appE-local-pfaffian-sign}
For the local rank-one type-II Pfaffian crossing
\((U_\delta, \mathrm{pf}_\delta)\) of
\S\ref{app:pfaffian-wall-rank-one},
\[
s_\delta^*\,\mathrm{pf}_\delta = -\mathrm{pf}_\delta,
\qquad
s_\delta^*\,(\mathrm{pf}_\delta^2) = \mathrm{pf}_\delta^2.
\]
Equivalently, the local Hall/Pfaffian sign is
\[
\epsilon_o^{\mathrm{loc}}(s_\delta) = (-1)^{N_\delta^{\mathrm{Pf}}} = -1,
\qquad
N_\delta^{\mathrm{Pf}} = 1.
\]
\end{proposition}

\begin{proof}
The fixed locus of \(s_\delta(t, u) = (t, -u)\) is \(\{u = 0\}\), and
the wall coincides with the tangent fixed locus, so the determinant
factor is split: \(\mathscr L^{\mathrm{tan}}_\delta\) on \(T_\delta\)
is the tangent factor and \(K^{\mathrm{cross}}_\delta\) is the normal
factor. The unit \(\upsilon_\delta\) is invariant by construction; the
skew matrix \(D_u\) has Pfaffian \(u\); and \(s_\delta\) sends
\(u \mapsto -u\) while fixing \(\upsilon_\delta\). Therefore
\[
s_\delta^*(\upsilon_\delta\,u) = \upsilon_\delta\,(-u) = -\upsilon_\delta\,u
= -\mathrm{pf}_\delta,
\]
and squaring gives the second equality. The Pfaffian rank exponent
\(N_\delta^{\mathrm{Pf}} = 1\) follows because the normal factor is the
single odd crossing \([\R \to \R]\). The displayed sign is
\((-1)^1 = -1\).
\end{proof}

\subsection{Match with the Maass multiplier}
\label{app:pfaffian-wall-maass-match}

The Maass multiplier \(\nu_{\Delta_5}\) of the weight-\(5\) Igusa cusp
form is the multiplier system of the section
\(\Delta_5 \in M_5(\Sp_4(\Z), \nu_{\Delta_5})\) of the automorphic
line bundle \(\mathcal L^5 \otimes \nu_{\Delta_5}\) on the
genus-two Siegel space \(\Sp(2,\Z) \backslash \Hpl_2\). Its values on
the type-II reflections in
\(\Orth(\La^{2,1})_+ = W^{(2)}(\La^{2,1})\) are computed in
\cite[Proposition 2.1]{GN} (citing Maass \cite[(1.3)--(1.5)]{MAASS:1})
via the exterior-square identification
\(\PSp_4(\Z) \simeq \Gamma_{3,2}\subset \mathrm P\Orth(\La^{3,2})_+\):

\begin{lemma}[Maass multiplier on simple type-II reflections]
\label{lem:appE-maass-on-type-ii}
The Maass character of \cite[Proposition 2.1]{GN}, in the
normalization of \cite[(1.3)--(1.5)]{MAASS:1}, evaluates on the
simple type-II reflections as
\[
\nu_{\Delta_5}(s_{\delta_i}) = -1, \qquad i = 1, 2, 3,
\]
and trivially on the three \(\mathrm{Aut}(\Poly_{II}) \simeq S_3\)
reflections \(s_{f_{-2} - f_2}, s_{f_2 - f_3}, s_{f_{-2} - f_3}\):
\[
\nu_{\Delta_5}(s_{f_{-2} - f_2})
= \nu_{\Delta_5}(s_{f_2 - f_3})
= \nu_{\Delta_5}(s_{f_{-2} - f_3}) = +1.
\]
The character is invariant under rescaling \(\Delta_5\) by a nonzero
constant: the theta-normalization
\([q^{1/2} r^{1/2} s^{1/2}] \Delta_5 = 64\) and the monic
\(D_5 = 64^{-1} \Delta_5\) give the same multiplier values on
generators.

\medskip
\emph{Derivation.}  Each \(s_{\delta_i}\) lifts under the exterior-square
isomorphism \(\PSp_4(\Z)\xrightarrow{\sim}\Gamma_{3,2}\) of
Theorem~\ref{thm:exceptional-iso} to a paramodular involution acting
on the genus-two Siegel space; \(\Delta_5\) transforms by the Maass
character on these involutions, valued \(-1\) on each type-II
generator (Maass \cite[(1.3)--(1.5)]{MAASS:1}, Gritsenko--Nikulin
\cite[Proposition 2.1]{GN}).  The trivial \(S_3\)-reflections
\(s_{f_{-2}-f_2}, s_{f_2-f_3}, s_{f_{-2}-f_3}\) lift to inner Weyl
permutations on the level-one paramodular and act with multiplier
\(+1\).
\end{lemma}

Combining Proposition \ref{prop:appE-local-pfaffian-sign} and Lemma
\ref{lem:appE-maass-on-type-ii}:

\begin{theorem}[Local orientation match on type-II generators]
\label{thm:appE-local-orientation-match}
Composing Proposition~\ref{prop:appE-local-pfaffian-sign} and
Lemma~\ref{lem:appE-maass-on-type-ii}, on the three simple type-II
reflections,
\[
\epsilon_o^{\mathrm{loc}}(s_{\delta_i}) = -1 = \nu_{\Delta_5}(s_{\delta_i}),
\qquad i = 1, 2, 3.
\]
The local Pfaffian sign and the Maass multiplier coincide on
generators of the type-II reflection group, relative to the retained
rank-one wall charts.
\end{theorem}

The match \(\epsilon_o^{\mathrm{loc}} = \nu_{\Delta_5}\) on the three
type-II simple reflections is the local seed of the global
orientation-character identity. Its extension to the whole
type-II Weyl group requires global Hall-stack transport of the
local crossing, the \((O2)\) obstruction.

\subsection{Extension to the type-II Weyl group}
\label{app:pfaffian-wall-W2-extension}

\subsubsection{Generation and a presentation}

The type-II Weyl group \(W^{(2)}(\La^{2,1}_{II})\) is generated by
\(s_{\delta_1}, s_{\delta_2}, s_{\delta_3}\) (Nikulin's reflection
theorem for \(\HypPlan \oplus \langle 2 \rangle\),
\cite[Theorem 4.1]{GN}). The generator pairings
\((\delta_i, \delta_j) = -2\) for \(i \ne j\) are obtuse and equal to
\(-2\); the corresponding Coxeter exponent
\(m_{ij}\) is the order of \(s_{\delta_i} s_{\delta_j}\), which is
infinite because the rotation \(s_{\delta_i} s_{\delta_j}\) acts on
the hyperbolic plane by a hyperbolic translation. Thus
\[
W^{(2)}(\La^{2,1}_{II})
= \langle s_{\delta_1}, s_{\delta_2}, s_{\delta_3} \mid
s_{\delta_i}^2 = 1 \rangle
\]
is the free product of three involutions, equivalently the
\((\infty, \infty, \infty)\)-triangle reflection group on the
hyperbolic plane.

\subsubsection{The determinant character}

The determinant character
\(\det : W^{(2)}(\La^{2,1}_{II}) \to \{\pm 1\}\) is the unique
homomorphism with \(\det(s_{\delta_i}) = -1\) on each generator. The relations \(s_{\delta_i}^2 = 1\) are respected since
\((-1)^2 = 1\), and there are no further relations, so \(\det\)
extends unambiguously. The Maass multiplier
\(\nu_{\Delta_5}|_{W^{(2)}(\La^{2,1}_{II})}\) is also a
\(\{\pm 1\}\)-character that equals \(-1\) on each generator (Lemma
\ref{lem:appE-maass-on-type-ii}); it therefore equals \(\det\):
\[
\nu_{\Delta_5}|_{W^{(2)}(\La^{2,1}_{II})}
= \det = (\text{sign character}).
\]
This is the global Maass character on the type-II Weyl group.

\subsubsection{The orientation character is a $\{\pm 1\}$-character}

The orientation character \(\epsilon_o\) on the Hall-stack moduli is,
by construction, a homomorphism
\(W^{(2)}(\La^{2,1}_{II}) \to \{\pm 1\}\) provided the orientation
cocycle on the half-Hilbert orbit
trivializes consistently with the local rank-one crossings of
\S\ref{app:pfaffian-wall-rank-one} across all generators. Under that
trivialisation, the global orientation character is determined by its
values on \(s_{\delta_1}, s_{\delta_2}, s_{\delta_3}\), each equal to
\(-1\) by Theorem \ref{thm:appE-local-orientation-match}. Hence
\[
\epsilon_o = \det = \nu_{\Delta_5}
\quad\text{on}\quad W^{(2)}(\La^{2,1}_{II}).
\]
This is the global orientation-character match relative to the retained
quotient orientation and the chosen Weyl transport.  Its proof in this
appendix establishes the \((O2)\) wall-atlas part: the trivialisation
of the orientation cocycle on the half-Hilbert orbit.

\subsection{The (O2) obstruction: half-Hilbert orbit transport}
\label{app:pfaffian-wall-O2}

\subsubsection{Statement of (O2)}

\begin{definition}[(O2) local Pfaffian wall normal form, global form]
\label{def:appE-O2}
The obstruction \((O2)\) is the assertion that, for every retained HN
height \(R\) and every retained type-II wall transport
\(s_\delta \in W^{(2)}(\La^{2,1}_{II})\), the local rank-one crossing
of \S\ref{app:pfaffian-wall-rank-one}, defined on the generic
wall chart \(U_{\delta, R}\), extends to a strictly compatible system
of charts on the half-Hilbert orbit
\(\mathfrak H^{\mathrm{half}}_R\), with the orientation cocycle on
this orbit trivialized by an explicit cochain compatible with the local
Pfaffian sign on each component, the boundary of each component, and
the overlap of every adjacent pair of components.
\end{definition}

\subsubsection{Local computation and global transport}

Theorem \ref{thm:appE-local-orientation-match} proves the local sign
identity on each of the three simple type-II generators by a single
chart computation: \emph{three local sign computations, one per
generator}. The global statement of
\S\ref{app:pfaffian-wall-W2-extension},
\(\epsilon_o = \nu_{\Delta_5}\) on the entire Weyl group, requires
that these three local seeds extend to a strictly coherent system on
the half-Hilbert orbit, namely the entire orbit of the half-Hilbert
moduli stack under the \(W^{(2)}\)-action. The component, boundary,
and overlap compatibility data of Definition \ref{def:appE-O2} are
\emph{not} produced by the local rank-one crossings alone.  They are
the content of the retained \((O2_{\mathrm{atlas}})\) datum transported
by Theorem~\ref{thm:G-O2-orbit-transport}.

\subsubsection{Two appearances of the constant $4096$}
\label{app:pfaffian-wall-O2-4096}

The integer \(4096\) appears in two structures, and the equality is
scalar only. On the orientation side, the
half-Hilbert orbit of the type-II reflection group on the
cosection-reduced self-Ext stack carries \(12\) binary signs, so the
orientation-cocycle support carries \(2^{12} = 4096\) sign assignments;
this is the reading on which the
\((O2_{\mathrm{atlas}})\) transport of Appendix~\ref{app:obstruction-discharges},
Theorem~\ref{thm:G-O2-orbit-transport}, operates. On the scalar side,
the Oberdieck--Pixton normalization
\[
Z^X_{\mathrm{OP}}(P_{\mathrm{OP}}, Q_{\mathrm{OP}}, T_{\mathrm{OP}})
= -\bigl(\chi_{10}^{\mathrm{OP}}\bigr)^{-1}
= -4096\, \Delta_5(Z)^{-2},
\qquad
\chi_{10}^{\mathrm{OP}} = D_5^2 = (64^{-1} \Delta_5)^2
= 4096^{-1} \Delta_5^2,
\]
exhibits \(4096 = 64^2 = ([q^{1/2} r^{1/2} s^{1/2}] \Delta_5)^2\) as the
squared theta-leading coefficient, with \(-4096 = (-1)\cdot 64^2\) the
OP scalar branch convention times the squared theta-leading
(\cite{OB,OPand}; arithmetic check by the product expansion).

The two values coincide because \(2^{12}=64^2\): the half-Hilbert orbit
sign count and the squared theta-leading are the same integer after
orientation has been forgotten.
Theorem~\ref{thm:G-R1-R2} of Appendix~\ref{app:obstruction-discharges}
makes the scalar separation precise: the orientation character
\(\epsilon_o = \nu_{\Delta_5}\) is fixed, relative to the retained
orientation and Weyl-transport data, by wall transport of the
sign-count, while the OP scalar magnitude \(4096\) enters
\(Z^X_{\mathrm{OP}} = -4096\,\Delta_5^{-2}\) through the squared
theta-leading.

\subsubsection{The retained (O2) wall-atlas route}

The retained \((O2_{\mathrm{atlas}})\) datum consists of:
\begin{enumerate}
\item[(O2.a)] A strictly compatible system of charts \(U_{\delta, R}\),
\(\delta \in W^{(2)}(\La^{2,1}_{II})\), on the half-Hilbert orbit at
height \(R\), each presenting the rank-one type-II Pfaffian crossing of
\S\ref{app:pfaffian-wall-rank-one}.
\item[(O2.b)] An orientation cocycle
\(c_o \in Z^2(W^{(2)}(\La^{2,1}_{II}); \mathbb F_2)\) on the orbit, with
explicit cochain trivialisation
\(\eta_o \in C^1(W^{(2)}(\La^{2,1}_{II}); \mathbb F_2)\) satisfying
\(\delta \eta_o = c_o\), compatible with the chosen charts.
\item[(O2.c)] Strict compatibility of the trivialisation with the
component decomposition, the boundary of each component, and the
overlap of every adjacent pair of components — equivalently, the
component, boundary, and overlap residual entries of the retained
wall-atlas datum vanish with explicit null-trivialisations.
\end{enumerate}
Once (O2.a)--(O2.c), the retained quotient orientation, and the chosen
Weyl transport are fixed, Theorem
\ref{thm:appE-local-orientation-match} extends to the global
Weyl-group character match
\(\epsilon_o = \nu_{\Delta_5}\) on \(W^{(2)}(\La^{2,1}_{II})\).

\subsection{First-timelike parity split as a recognition obligation}
\label{app:pfaffian-wall-recognition}

The orientation-character match of \S\ref{app:pfaffian-wall-W2-extension}
is the conditional \emph{multiplier} comparison. The recognition
theorem \(\mathcal P_X \cong \fg_{\Delta_5}\) of Chapter
\ref{ch:pfaffian-dirac} requires further data, namely full
parity dimensions per root degree, not just signed dimensions. The
first nontrivial test occurs at the imaginary simple root
\(\delta_{123} = \delta_1 + \delta_2 + \delta_3\):

\begin{lemma}[First-timelike parity split]
\label{lem:appE-first-timelike-parity}
The signed root multiplicity at \(\delta_{123}\) is
\[
\smult\,\fg_{\Delta_5, \alpha(\delta_{123})} = f(1, 1) = -64,
\]
and the full target parity computed from the Gritsenko--Nikulin
presentation is
\[
\bigl(\rootmult_{\overline 0}(\delta_{123}),\,
      \rootmult_{\overline 1}(\delta_{123})\bigr)
= (29, 93),
\qquad
29 - 93 = -64.
\]
The decomposition of the even part is
\[
29 = 2 + 3 \cdot 9,
\]
where \(2 = \tfrac{1}{3}\,\mu_{\mathrm{Mob}}(1)\,3!\) is the Witt dimension of the
free Lie algebra on three generators in multidegree \((1, 1, 1)\) and
\(3 \cdot 9\) is the three real--isotropic mixings each contributing the
eta-9 multiplicity \(\tau = 9\) of the doubled isotropic ray. The odd
part \(93\) is the absolute value of the timelike Borcherds correction
\(m(\delta_{123})\), with
\[
m(\delta_{123}) =
-\bigl[\,q^1 r^1 s^1\,\bigr]\bigl((qrs)^{-1/2}64^{-1}\Delta_5\bigr)
= -93.
\]
\end{lemma}

\begin{proof}
The signed multiplicity is \(f(1, 1) = -64\) by direct expansion of
\(\phi_{0,1}\) in the Eichler--Zagier normalization
\cite[\S2.2]{EZ:1}. The full parity decomposition
is the Gritsenko--Nikulin presentation reading: the even part counts
free Lie words in three real-root letters, the odd part counts simple
imaginary timelike roots. Witt's dimension formula
\(W_{\mathrm{multidegree}}(1,1,1)
= \tfrac{1}{n} \sum_{d \mid \gcd} \mu_{\mathrm{Mob}}(d)\,(n/d)! / \prod (n_i/d)!\) with
\(n = 3\), \(\gcd(1,1,1) = 1\) gives \(\tfrac{1}{3} \cdot 1 \cdot 6 = 2\).
The three real--isotropic mixings each contribute
\(\tau = 9\) on the primitive isotropic rays (the eta-9 expansion
\(\prod_{k \ge 1}(1 - q^k)^9\) of \cite[Theorem 2.1, (5.1)]{GNII}); summing over
the three independent isotropic directions gives \(3 \cdot 9 = 27\).
Hence \(29 = 2 + 27\). The odd part is \(-m(\delta_{123}) = 93\) by the
Gritsenko--Nikulin imaginary simple-root formula.
\end{proof}

\begin{corollary}[Signed dimension is not recognition]
\label{cor:appE-signed-not-recognition}
A graded super vector space with parity dimensions
\((\dim_{\overline 0}, \dim_{\overline 1}) = (29, 93)\) and zero
bracket has the same signed dimension \(-64\) as the
\(\fg_{\Delta_5, \alpha(\delta_{123})}\) component but the wrong Lie
superalgebra structure. Recognition of \(\mathcal P_X\) with
\(\fg_{\Delta_5}\) at degree \(\delta_{123}\) requires both parity
dimensions \(29\) and \(93\) and the Borcherds--Kac--Moody bracket; the
determinant identity alone does not imply this recognition statement.
\end{corollary}

\subsubsection{Higher-degree consistency}

The same Borcherds--Kac arithmetic gives higher-degree signed checks
and the finite rows where full target parity is known:

\begin{itemize}
\item \emph{Doubled isotropic.} For \(\delta = 2 a_{ij}\), \(i \ne j\),
\[
(\smult,\, \tau,\, \text{gap}) = (10,\, 9,\, 1),
\]
the signed multiplicity \(10\) splits as eta-9 multiplicity \(9\) plus
the Kac null-root residual \(1\) of the rank-two affine subroot
system, for each of the three pairs \((i,j)\).

\item \emph{Height-four timelike.} For the three rows
\[
C_{k,2}=a_{ij}+2\delta_k,\qquad \{i,j,k\}=\{1,2,3\},
\]
namely coefficient vectors \((2,1,1),(1,2,1),(1,1,2)\),
\[
(\smult,\, m,\, \text{gap}) = (108,\, 90,\, 18),
\]
each with the same triple of integers.  These rows are not Weyl
translates of \(\delta_{123}\); rather
\[
s_{\delta_k}(\delta_{123})=a_{ij}+3\delta_k=C_{k,3}.
\]
Thus \(C_{k,2}\) is a signed target row.  The integer \(18=108-90\)
is the signed lower-bracket residual at height four, not a full parity
block.

\item \emph{Real-string exponents.} The Borcherds--Kac real-root strings
through \(\delta_{ij}\)-style mixed degrees give
\((\text{real-real exponent}) = 3\) and
\((\text{complement-isotropic exponent}) = 5\); the timelike string
through \(\delta_{123}\) has exponent \(3\). All three are pairwise
consistent under Weyl reflection.
\end{itemize}
These integers are the degree-wise Borcherds--Kac arithmetic used in
the recognition criterion.

\subsection{Local Pfaffian sign computation}
\label{app:pfaffian-wall-sign-computation}

Let \(\delta\in\{\delta_1,\delta_2,\delta_3\}\subset\La^{2,1}\), with
\((f_2,f_{-2})=-1\) and \((f_3,f_3)=2\).  On a normal chart
\[
U_\delta=T_\delta\times(-\epsilon,\epsilon),\qquad
s_\delta(t,u)=(t,-u),\qquad
\Hpl_\delta\cap U_\delta=T_\delta\times\{0\},
\]
choose an oriented tangent Pfaffian unit \(\upsilon_\delta\) satisfying
\(s_\delta^*\upsilon_\delta=\upsilon_\delta\).  The rank-one normal
complex is
\[
K^{\mathrm{nor}}_\delta=[\mathbb R\xrightarrow{u}\mathbb R],
\qquad
D_u=
\begin{pmatrix}
0&u\\
-u&0
\end{pmatrix},
\qquad
\operatorname{Pf}(D_u)=u .
\]
Thus
\[
\operatorname{pf}_\delta=\upsilon_\delta u,\qquad
s_\delta^*\operatorname{pf}_\delta=-\operatorname{pf}_\delta,\qquad
s_\delta^*\operatorname{pf}_\delta^2=\operatorname{pf}_\delta^2 .
\]
The normal Pfaffian exponent is \(N_\delta^{\mathrm{Pf}}=1\), hence
\[
\epsilon_o^{\mathrm{loc}}(s_\delta)
=(-1)^{N_\delta^{\mathrm{Pf}}}
=-1
=\nu_{\Delta_5}(s_\delta)
\]
for the three simple type-II reflections, by the Maass-character
calculation of Gritsenko--Nikulin \cite[Proposition~2.1]{GN}.

The \((O2_{\mathrm{atlas}})\)-transport formulas of
Appendix~\ref{app:eight-finite-constructions}
extends this local computation to the global half-Hilbert orbit by
component, boundary, and overlap compatibility.

\subsection{O2 relative summary}
\label{app:pfaffian-wall-open-obligations}

\begin{itemize}
\item \textbf{(O2.a) Strictly compatible chart system on the half-Hilbert
orbit.} The local rank-one crossing is fixed at each simple
type-II generator (\S\ref{app:pfaffian-wall-rank-one}); strict
compatibility across all elements of \(W^{(2)}(\La^{2,1}_{II})\) is
the half-Hilbert orbit chart system retained in
\((O2_{\mathrm{atlas}})\) and transported by
Theorem~\ref{thm:G-O2-orbit-transport}.
\item \textbf{(O2.b) Orientation cocycle trivialisation.}
\(c_o \in Z^2(W^{(2)}(\La^{2,1}_{II}); \mathbb F_2)\) is the
orientation cocycle; its trivializing cochain
\(\eta_o \in C^1(W^{(2)}(\La^{2,1}_{II}); \mathbb F_2)\) is part of
the chosen Weyl-transport datum of Theorem~\ref{thm:G-O1-plus} and
is transported in
Theorem~\ref{thm:G-O2-orbit-transport}.
\item \textbf{(O2.c) Component, boundary, overlap compatibility.}
The residual entries
\(\mathfrak o^{\mathrm{chart}}_{\delta, R}\),
\(\mathfrak o^{\mathrm{wall}}_{\delta, R}\),
\(\mathfrak o^{\mathrm{split}}_{\delta, R}\),
\(\mathfrak o^{\mathrm{unit}}_{\delta, R}\),
\(\mathfrak o^{\mathrm{rank}}_{\delta, R}\),
\(\mathfrak o^{\mathrm{tr}}_{\delta, R^+ R}\) of
Proposition \ref{prop:rank-one-crossing-o2-sign} are fixed at finite
height \(R\) for the local rank-one models; the global compatibility
across the half-Hilbert orbit is retained in
\((O2_{\mathrm{atlas}})\) and transported by
Theorem~\ref{thm:G-O2-orbit-transport}.
\item \textbf{\((R1)\)--\((R2)\) scalar separation.} Separated in
Appendix~\ref{app:obstruction-discharges},
Theorem~\ref{thm:G-R1-R2}: the orientation-cocycle support on the
half-Hilbert orbit has size \(2^{12}\) (R1), and this integer is equal
to \(64^2\), the theta-leading square on the scalar side (R2).  This is
not an identification of the two data.
\end{itemize}

Relative to the chosen Weyl lifts and \((O2_{\mathrm{atlas}})\),
Theorems~\ref{thm:G-O1-plus}, \ref{thm:G-O2-orbit-transport}, and
\ref{thm:G-R1-R2} extend Theorem
\ref{thm:appE-local-orientation-match} from the three local generators
to the global character match
\(\epsilon_o = \nu_{\Delta_5}\) on \(W^{(2)}(\La^{2,1}_{II})\).
